\chapter*{General Introduction}
\addcontentsline{toc}{chapter}{General Introduction}
\markboth{General Introduction}{General Introduction}

A consulting firm earns its margin one project at a time. A contract fixes the price and the
number of days sold; whether any money is left at the end depends on what the work actually costs
to deliver, and on when the client is billed for it. Those three figures move independently of
each other. Lose sight of one, and the news that a project lost money arrives after it has closed,
which is the one moment when nothing can be done about it.

At ST2i that arithmetic lives in a spreadsheet. The file is called a project review sheet: some
twenty linked tabs, more than two thousand formulas, one copy per project. It holds a great deal.
Cost rates for every resource and every year, the workload planned month by month, billing
milestones written as percentages of the contract, mission expenses, the governance registers, and
a complete earned-value engine that computes the margin at the end of it. Nobody ever wrote the
company rules down as a specification. This file is the specification.

It is also breakable. One file per project means there is no single copy anyone can point to as
the true one. Formulas that reach across tabs fail quietly, and the workbook handed over at the
start of this project has cells in error sitting in it. Seeing the whole portfolio means opening
every file in turn. And nothing in a spreadsheet stops a developer from opening the tab that shows
the margin on the project they are working on.

\textbf{PMS}, for \textit{Project Management System}, is the platform built to take that work
over. Two things were settled before any code was written. It had to reproduce the financial model
as the workbook computes it, down to the cost rates, the planned and consumed workload, the
estimate at completion, the margin and the billing progress, rather than a tidier approximation of
it: a figure that disagrees with the workbook is a figure nobody at ST2i will trust. And the
question of who may see what had to be answered in the database rather than in the code, because
those rules change, and a release is too high a price for granting somebody a right.

It was built in ten sprints grouped into four releases, each one demonstrable on its own before
the next began. This report follows that order.

\textbf{Chapter~1} is about where the project comes from: the company, the way its projects are
steered today, the workbook read sheet by sheet, and what all of that adds up to as a problem
worth solving.

\textbf{Chapter~2} turns the problem into requirements. Who the users are and what each of them
may reach, what the platform has to do and how well it has to do it, the rules the business
imposes on it, and the backlog all of that became.

\textbf{Chapter~3} is the design: the architecture chosen and the reasoning behind it, the
decisions that apply across every module, and the database the whole thing stands on.

\textbf{Chapters~4 to~7} are the construction, a release to a chapter. First the foundation and
its security, then the day-to-day management of projects and workload, then the financial engine
the project exists for, and last the work of turning a finished application into a deliverable
one.

The \textbf{general conclusion} looks back at what was asked for at the start, says how much of it
arrived, and sets out what could follow.
